
\documentclass[../../../physics-standard_questions_all.tex]{subfiles}


\begin{document}

図のように，真空中において，同形の3枚の金属極板X，Y，Zを
平行に狭い間隔で向き合わせてあり，Zはアース（接地）されている．
XとYの間隔は$d$，YとZの間隔は$2d$である．極板の面積を$S$，
真空の誘電率を$\varepsilon_0$とし，端の効果は無視する．
はじめ，いずれの極板も帯電していない．

\begin{enumerate}

    \item Xに電荷$Q_0$を与え，安定するまで待った．
    このとき，X下面の電荷$x_1$，Y下面の電荷$y_1$を求めよ．

    \item 一旦はじめの状態に戻し，電源によりXの電位を$V_0$に設定して，安定するまで待った．
    このとき，X下面の電荷$x_2$，Y下面の電荷$y_2$を求めよ．

    \item 一旦はじめの状態に戻し，Yに電荷$Q_0$を与え，
    XとZを導線でつないで，安定するまで待った．
    このとき，X下面の電荷$x_3$，Y下面の電荷$y_3$を求めよ．

\end{enumerate}

\vspace{0.2\baselineskip}
{\centering
    \includegraphics{\subfix{fig/fig_em_cb_05/fig_em_cb_05_00.pdf}}
\par}
\vspace{0.2\baselineskip}

\end{document}



